% Source: https://artofproblemsolving.com/community/c3h11311_mathcounts_problems_3
% Stable identity: aops-topic-11311-q04
Reflecting in the horizontal line $y=-2$ sends a point $(x,y)$ to $(x,-4-y)$ (the two points have the same $x$ and their $y$-values average to $-2$). So the image of $y=x^2+3$ is
\[
y=-4-(x^2+3)=-x^2-7 .
\]

\begin{asy}
usepackage("amsmath");
size(5cm);
defaultpen(fontsize(10pt));
import graph;
real f(real x){return x^2+3;} real g(real x){return -x^2-7;}
draw((-4,0)--(4,0),gray,Arrow(TeXHead));draw((0,-11.5)--(0,7.5),gray,Arrow(TeXHead));
label("$x$",(4,0),E,gray);label("$y$",(0,7.5),N,gray);
draw(graph(f,-2,2),blue+linewidth(0.9));draw(graph(g,-2,2),red+linewidth(0.9));
draw((-3.5,-2)--(3.5,-2),dashed);label("$y=-2$",(3.5,-2),E);
dot((0,3),blue);dot((0,-7),red);
label("$(0,3)$",(0,3),SE,blue);label("$(0,-7)$",(0,-7),NE,red);
label("$y=x^2+3$",(2,7),E,blue);label("$y=-x^2-7$",(2,-11),E,red);
label("\mbox{}",(min(currentpicture,user=true)+max(currentpicture,user=true))/2,black); // keeps the text after the figure black when LaTeX typesets the labels inline
\end{asy}

Since $-x^2\le0$, the maximum value is $\boxed{-7}$, at the vertex $(0,-7)$, the image of the original vertex $(0,3)$.
